\chapter{Related Work}
\label{ch:related}

This chapter reviews what is known about drag- and lift-based swimming in animals and on foils, how swimming
legged robots have been designed and simulated, and which gap this thesis addresses.

\section{Drag- and Lift-Based Swimming in Animals}
\label{sec:rw-bio}

Fish \cite{Fish1996} classified mammalian swimming by the force that produces thrust. Paddlers such as the
muskrat, mink or dog move a foot or hand backwards against the water in a power stroke and return it in a
recovery stroke. Their propulsive efficiency is at most about 0.33. Swimmers that oscillate a fluke or flipper
across the swimming direction generate lift and reach efficiencies above 0.8. Fish argued that the transition
from paddling to lift-based swimming in the evolution of aquatic mammals was driven by this efficiency
difference and that intermediate forms combine both principles.

The muskrat is the best-documented paddler \cite{Fish1984}. Its power stroke lasts \SI{0.18}{\second} and
its recovery \SI{0.22}{\second}, a duty factor of 0.45. During recovery it folds the toes together, reducing
the frontal area of the foot by 55.5\,\%, yet the recovery drag still amounts to about one third of the
power-stroke thrust. Its maximum mechanical efficiency is 0.33. Dogs use a stereotyped diagonal paddling
gait, in which each foot moves on a closed path and the fore- and hindlimbs of opposite sides move together
\cite{Fish2021DogPaddle}. Ducks paddle with webbed feet whose hydrodynamic performance resembles that of oars
rather than delta wings \cite{Ribak2023}.

Cetaceans are the reference for lift-based swimming. Bottlenose dolphins swim with a propulsive efficiency
of about 0.8 \cite{Fish1993}. Odontocetes cruise at Strouhal numbers of $0.26\pm0.05$ \cite{Rohr2004}, and
flying and swimming animals in general cruise at $0.2<\mathit{St}<0.4$, the range of high power efficiency of
oscillating foils \cite{Taylor2003}. Walker and Westneat \cite{Walker2000} compared rowing and flapping
appendages with a blade-element model and concluded that rowing is favourable at low speeds and for
manoeuvres and acceleration, whereas flapping is more efficient at higher speeds. For BODY2 this suggests the division of
labour tested in this thesis, a paddle for starting and low speed and a fluke for cruising, and makes the speed at which one
overtakes the other the quantity to be found.

\section{Flapping-Foil Hydrodynamics}
\label{sec:rw-foil}

A foil that heaves and pitches with a phase difference produces thrust through a leading-edge vortex on each
half stroke and a reverse von K\'arm\'an street in its wake \cite{Triantafyllou1991}. Wake stability analysis
predicts the highest efficiency at $0.25\lesssim\mathit{St}\lesssim0.35$ \cite{Triantafyllou1991}. Anderson et
al.\ \cite{Anderson1998} measured efficiencies up to 0.87 for a NACA\,0012 foil at a heave amplitude of
$0.75$ chord, a maximum angle of attack of about \SI{20}{\degree} and a pitch lead of \SI{75}{\degree}. Read
et al.\ \cite{Read2003} extended these measurements and found a broad efficiency plateau of 0.5--0.6 for
maximum angles of attack between \SI{15}{\degree} and \SI{25}{\degree} and the largest thrust coefficients,
around 2.4, at $\mathit{St}\approx0.6$ and angles of attack of \SIrange{30}{35}{\degree}, where efficiency drops
to below 0.5. Hover et al.\ \cite{Hover2004} showed that the time history of the angle of attack matters as
much as its maximum: an angle of attack held nearly constant over the stroke extends the efficient range to
higher Strouhal numbers and thrust. Schouveiler et al.\ \cite{Schouveiler2005} reported thrust coefficients
and efficiencies with error bars over a range of Strouhal numbers and angles of attack; one of their points
is used as the validation case in \cref{sec:ver-v1}. Floryan et al.\ \cite{Floryan2017} derived scaling laws
in which thrust grows with the square of the Strouhal number and efficiency favours large amplitudes at low
frequencies. For a self-propelled swimmer, the Strouhal number is not a free parameter but follows from the
balance of thrust and drag \cite{Eloy2012}. Three-dimensional effects reduce thrust at low aspect ratio;
cetacean flukes have aspect ratios of about 2--6 \cite{Ayancik2020}. For the fluke of BODY2 these results fix the target: an
angle of attack of about \SIrange{15}{25}{\degree} and a Strouhal number of 0.2--0.4, which its pitch must maintain while the
swimming speed changes.

\section{Paddling Mechanics}
\label{sec:rw-paddle}

A paddle produces net thrust only through an asymmetry between power and recovery stroke, either spatial
(a smaller area during recovery) or temporal (a slower recovery) \cite{HerreraAmaya2024}. Kim and Gharib
\cite{Kim2011} measured the flow around a flat plate performing a paddling rotation and showed that the thrust
follows the growth and shedding of the starting vortex rather than the quasi-steady square of the plate
velocity. Crustacean-inspired propulsors reach their best performance when the appendage opens quickly at
the start of the power stroke and closes early in the recovery \cite{Santos2026}. A tendon coupling that
shortens the transition between the open and folded states of a webbed robotic foot doubled its propulsive
efficiency \cite{Lin2024}. For beaver-like robots, the hydrodynamics of bendable webbed feet have been modelled
with quasi-steady methods and measured in flumes \cite{Chen2021,Chen2022}. For the fan paddle of BODY2 the levers are therefore
the folded area, the timing of the fold and the velocity law of the power stroke; all three are varied in \cref{sec:res-drag}.

\section{Swimming Legged Robots}
\label{sec:rw-robots}

Picardi et al.\ \cite{Picardi2023} reviewed underwater legged robots; most walk on the bottom and few swim.
Li et al.\ \cite{Li2019} analysed the kinematics and hydrodynamics of a robotic dog paddling with rigid
limbs. Qu et al.\ \cite{Qu2025} built an amphibious robotic dog and measured single-leg and whole-robot
thrust for paddling gaits with different power-stroke fractions; a shorter power stroke produced more thrust
but also larger roll. Wang et al.\ \cite{Wang2025} simulated quadrupedal paddling with an immersed-boundary
CFD solver. They varied the initial swing angle, the swing amplitude and the power-stroke fraction of rigid,
non-folding legs at zero forward speed, and found that the thrust of four legs exceeds four times the
single-leg thrust by up to 87\,\%. Han et al.\ \cite{Han2025} trained an LSTM model of the leg forces from
tank measurements and used it for multi-objective gait optimization. Yao et al.\ \cite{Yao2023} instead added
thrusters to a quadruped and learned a propulsion controller.

Lift-based propulsion has been demonstrated on robots with dedicated flippers. Baines et al.\
\cite{Baines2022} built a turtle robot with morphing limbs that can paddle and flap with the same limb; the
flapping gait produced thrust over most of the stroke and had the lower cost of transport, while paddling
gave higher accelerations. The DuckyDog robot uses variable-stiffness legs with passive fins
\cite{Wang2025DuckyDog}. Within the same chair as this thesis, Heber \cite{Heber2026} optimized drag-based
paddling gaits of an amphibious quadruped by direct collocation over a reduced-order hydrodynamic model
calibrated against an SPH simulation. \Cref{tab:rw-robots} summarizes these studies. Most measure or simulate thrust at rest, and
none compares a drag-based and a lift-based propulsor on the same leg across forward speed, which is what the choice of propulsor
for BODY2 requires.

\begin{table}[htb]
  \centering
  \caption[Swimming legged robots]{Swimming legged robots and this thesis. Speeds as reported; $^\dagger$derived from the reported course
    length and time; -- not reported.}
  \label{tab:rw-robots}
  \footnotesize
  \begin{tabularx}{\textwidth}{@{}>{\raggedright\arraybackslash}p{30mm}>{\raggedright\arraybackslash}p{34mm}>{\raggedright\arraybackslash}p{30mm}>{\raggedright\arraybackslash}X>{\raggedright\arraybackslash}p{24mm}@{}}
    \toprule
    Study & Propulsor & Method & Thrust vs.\ forward speed & Swimming speed \\
    \midrule
    Li et al.\ \cite{Li2019} & rigid paddling legs (pneumatic soft actuators) & kinematic and drag analysis, prototype & no & -- \\
    Qu et al.\ \cite{Qu2025} & rigid two-segment paddling legs & force measurement at rest, free swimming & no & \SI{0.16}{\metre\per\second} (0.54\,BL/s) \\
    Wang et al.\ \cite{Wang2025} & rigid paddling legs & immersed-boundary CFD, tethered & no & -- \\
    Han et al.\ \cite{Han2025} & paddling legs with web plates & tank measurements, LSTM model, gait optimization & yes, \SIrange{0}{0.3}{\metre\per\second} & $\approx$\SI{0.1}{\metre\per\second}$^\dagger$ \\
    Baines et al.\ \cite{Baines2022} & morphing limb: paddling and flapping & force measurement at rest, CFD at one speed & no & -- \\
    DuckyDog \cite{Wang2025DuckyDog} & legs with passive fins & free swimming & no & 0.30\,BL/s \\
    Heber \cite{Heber2026} & rigid paddling legs & reduced-order model, direct collocation & no & -- \\
    \midrule
    This thesis & folding fan paddle and fluke on the same leg & reduced-order optimization, overset CFD, hardware test & yes, \SIrange{0}{0.4}{\metre\per\second} & CFD \SIrange{0.26}{0.30}{\metre\per\second} (1.1--1.3\,BL/s); hardware \SI{0.14}{} / \SI{0.26}{\metre\per\second} mean from rest at 2 / 4\,Hz \\
    \bottomrule
  \end{tabularx}
\end{table}


\section{Simulation Approaches for Swimming Robots}
\label{sec:rw-sim}

The simulation methods span several orders of magnitude in cost. Reduced-order models apply quasi-steady drag,
added mass and buoyancy to each link \cite{Morison1950,Fossen1994}. They are fast enough for optimization
\cite{Chen2021,Heber2026,Lee2025}, but their coefficients must be identified, and a plain drag model contains
no lift. Particle methods such as SPH and grid-based CFD resolve the flow, but one simulated second takes
minutes to hours. Overset (Chimera) grids are an established way to simulate large prescribed body motions in
OpenFOAM \cite{Weller1998,Windt2020}; the immersed-boundary method is an alternative \cite{Wang2025}.
Lattice-Boltzmann solvers on adaptive block grids, such as the commercial Flare Live \cite{FlareLive}, trade
resolution for real-time speed and can simulate a whole robot with its joints and actuators. For this thesis the consequence is
the division of tools described in \cref{ch:intro}: the reduced-order model for the search and the CFD for the numbers.

\section{Research Gap}
\label{sec:rw-gap}

The literature establishes that lift-based flapping is more efficient than drag-based paddling in animals
and that paddling can be improved by folding, timing and velocity shaping. Robotic studies of legged swimming
have concentrated on drag-based strokes of rigid or folding feet, mostly at zero forward speed
\cite{Wang2025,Qu2025}, and on reduced-order models whose lift content is either absent or uncalibrated.
To the author's knowledge, no study has compared an optimized drag-type paddle and a lift-type fluke on the
same legged-robot leg across forward speed with a resolved flow simulation, identified what limits each
propulsor, and translated the result into swimming speeds of the whole robot. The following combination is
therefore missing and is addressed here:
\begin{itemize}
  \item a reduced-order optimization that includes actuator speed limits and a lift term, to see which
        principle an optimizer chooses;
  \item single-leg CFD of the same robot leg with both propulsors, each with an optimized motion, from rest to
        cruising speed;
  \item a self-propulsion estimate and a cross-check with a whole-robot real-time solver.
\end{itemize}
